\documentclass[12pt,a4paper]{article}
\usepackage[T1]{fontenc} 
\usepackage[utf8]{inputenc}
\usepackage{amsmath}
\usepackage{amsthm}
\usepackage{amssymb}
\usepackage{graphicx}
\usepackage{tabularx}
\usepackage{multirow}
\usepackage{float}
\usepackage{listings}
\usepackage[hidelinks]{hyperref}
\usepackage[polish]{babel}
\usepackage[useregional=numeric]{datetime2}
\usepackage[a4paper,margin=1in]{geometry}

\newtheorem{lem}{Lemat}
\newtheorem{thm}{Twierdzenie}
\theoremstyle{definition}
\newtheorem{defn}{Definicja}

\DeclareMathOperator{\ability}{ability}
\DeclareMathOperator{\need}{need}
\DeclareMathOperator{\assign}{assign}
\DeclareMathOperator{\missing}{missing}

\renewcommand{\lstlistingname}{Przykład}
\lstset{
	captionpos=b
}

\begin{document}

\begin{titlepage}
	\centering
	{\Large Wydział Matematyki i Nauk Informacyjnych\\Politechniki Warszawskiej \par}
	\vspace{1cm}
	\includegraphics[width=0.2\textwidth]{res/img/logo.png} \par
	\vspace{4cm}
	{\LARGE Teoria Algorytmów i Obliczeń\\Laboratorium - Etap 1 \par}
	\vspace{0.5cm}
	{\LARGE Specyfikacja wstępna \par}
	\vspace{2cm}
	{\large Adrian Bednarz,\\Bartłomiej Dach,\\Tymon Felski \par}
	\vspace{2cm}
	{\large Wersja 1.0 \par}
	\vspace{0.5cm}
	{\large \today \par}
\end{titlepage}

\noindent
Lista zmian:\\[.5\baselineskip]
\begin{tabularx}{\textwidth}{|l|l|X|l|}
	\hline
	\textbf{Data} & \textbf{Autor} & \textbf{Opis} & \textbf{Wersja} \\
	\hline
	15.10.2017 & Tymon Felski & Stworzenie szablonu dokumentu & 1.0 \\
	\hline
\end{tabularx}

\newpage
\tableofcontents
\newpage

\section{Opis problemu}
\label{sec:description}
Niniejszy rozdział poświęcony jest dokładnemu opisaniu podstawowej wersji zadanego problemu.\\

\noindent
Dane są zbiory:
\begin{itemize}
	\item $E$ - \textbf{ekspertów} realizujących projekty,
	\item $S$ - \textbf{umiejętności} posiadanych przez ekspertów,
	\item $P$ - \textbf{projektów} do zrealizowania.\\
\end{itemize}

\noindent
Każdemu ekspertowi przypisany jest wektor binarny opisujący posiadane przez niego umiejętności. Przykładowo, jeżeli ekspert posiada umiejętność $i$, to w wektorze umiejętności odpowiadającemu temu ekspertowi na $i$-tym miejscu znajduje się znak $1$, w przeciwnym wypadku - $0$.

\begin{verbatim}
Załóżmy, że liczność zbioru umiejętności S jest równa 5. Wówczas ponume-
rujemy umiejętności rozważane w problemie liczbami z zakresu [1, 5].

Pewien ekspert ze zbioru E posiada umiejętności 2, 3 i 5. Wówczas wektor
opisujący jego umiejętności to:

                            [0, 1, 1, 0, 1]
\end{verbatim}
\begin{lstlisting}[caption={Wektor opisujący umiejętności eksperta}]
\end{lstlisting}

\noindent
Każdemu projektowi przypisany jest wektor liczbowy opisujący zapotrzebowanie na ekspertów posiadających dane umiejętności. Przykładowo, jeżeli do realizacji projektu potrzeba trzech ekspertów posiadających umiejętność $i$, to w wektorze umiejętności odpowiadającemu temu projektowi na $i$-tym miejscu znajduje się liczba $3$.

\begin{verbatim}
Utrzymując założenie o liczności zbioru umiejętności z poprzedniego przy-
kładu, rozważmy pewien projekt ze zbioru P. Niech jego zapotrzebowanie na
ekspertów posiadających umiejętności 1 i 4 wynosi odpowiednio 4 i 3, a na
pozostałe - 0. Wówczas wektor opisujący zapotrzebowanie tego projektu to:
 
                            [4, 0, 0, 3, 0]
\end{verbatim}
\begin{lstlisting}[caption={Wektor opisujący zapotrzebowanie projektu}]
\end{lstlisting}

\noindent
Wszystkie projekty realizowane są w tym samym oknie czasowym, tzn. prace nad każdym z nich rozpoczynają się w momencie $t_0$ i kończą w późniejszym momencie $t_k$. Oznacza to, że jeżeli dany ekspert zostanie zatrudniony do pracy nad projektem $P_1$, to nie będzie mógł brać udziału w równoległym projekcie $P_2$. Ponadto każdy ekspert podczas pracy nad projektem może wykorzystywać tylko jedną z posiadanych umiejętności i nie może jej zmienić w trakcie trwania prac.\\

\noindent
Prace nad danym projektem zostaną zakończone nawet jeżeli nie zostanie mu przydzielona wymagana liczba ekspertów posiadających potrzebne umiejętności, określona przez wektor liczbowy odpowiadający temu projektowi. Będzie on natomiast zrealizowany gorzej niż w przypadku, gdyby przypisana została odpowiednia liczba ekspertów. Może się również zdarzyć, że najbardziej optymalne okaże się takie przypisanie ekspertów, że nad pewnym projektem nie będzie pracował nikt.\\

\noindent
Wówczas mamy do czynienia z brakami. Poprzez braki rozumiemy różnicę pomiędzy zapotrzebowaniem projektu na ekspertów o danych umiejętnościach a rzeczywistym przydziałem. Aby wyznaczyć braki w danym projekcie, należy odjąć wektor zapotrzebowania projektu na ekspertów od wektora zawierającego informację o ekspertach przydzielonych do tego projektu i zsumować elementy uzyskanej różnicy. Dokładna definicja tego pojęcia znajduje się w rodziale zawięrającym dowód poprawności (definicja nr 8).

\begin{verbatim}
Niech wektorem opisującym zapotrzebowanie projektu na ekspertów będzie:
                            [4, 0, 0, 3, 0]
Załóżmy, że do tego projektu zostali przypisani eksperci opisani przez:
                            [1, 0, 1, 0, 1] (wykorzystuje 1)
                            [1, 0, 0, 0, 1] (wykorzystuje 1)
                            [1, 1, 0, 1, 0] (wykorzystuje 4)
                            [0, 0, 0, 1, 1] (wykorzystuje 4)
Stąd przydział ekspertów do tego projektu można opisać wektorem:
                            [2, 0, 0, 2, 0]
Braki w tym projekcie obliczymy następująco:
    sum([4, 0, 0, 3, 0] - [2, 0, 0, 2, 0]) = sum([2, 0, 0, 1, 0]) = 3
\end{verbatim}
\begin{lstlisting}[caption={Wektor opisujący zapotrzebowanie projektu}]
\end{lstlisting}

\noindent
Naszym celem jest zminimalizowanie braków w obrębie wszystkich projektów (sumy wszystkich braków), czyli znalezienie optymalnego przydziału ekspertów do projektów.

\section{Problem znajdowania maksymalnego przepływu}
\label{sec:maxflow}
Okazuje się, że problem opisany w~sekcji~\ref{sec:description} można uogólnić
do~znanego problemu znajdowania maksymalnego przepływu w~sieciach. W~tej sekcji
zdefiniowane są podstawowe pojęcia potrzebne do~opisu tego problemu.

\begin{defn}
\textbf{Siecią} nazywamy czwórkę uporządkowaną $S = (G,c,s,t)$, gdzie:
\begin{itemize}
	\item $G = (V,E)$ jest grafem skierowanym,
	\item $c : E \to \mathbb{N}$ to tzw. funkcja przepustowości,
	\item $s,t \in V, s \neq t$ są dwoma wyróżnionymi wierzchołkami grafu $G$
	--- kolejno źródłem i ujściem sieci.
\end{itemize}
\end{defn}

\begin{defn}
\textbf{Przepływem} w~sieci $S$ nazywamy funkcję $f : E \to \mathbb{N}$
spełniającą następujące warunki:
\begin{enumerate}
	\item $\forall_{e \in E} \quad 0 \leqslant f(e) \leqslant c(e)$,
	\item $\displaystyle(\forall v \in V - \{ s,t \})
	\sum_{u : \; uv \in E} f(uv) = \sum_{u : \; vu \in E} f(vu)$ 
	--- tzw. prawo Kirchhoffa.
\end{enumerate}
\end{defn}

\noindent
W~ogólniejszym przypadku funkcje przepustowości i~przepływu mogą mieć wartości
nieujemne rzeczywiste, lecz założenie o~całkowitości zapewnia, że~algorytmy
wyznaczające maksymalny przepływ zawsze zakończą działanie.\\

\noindent
Prawo Kirchhoffa stanowi, że~suma wartości przepływu na~krawędziach wchodzących
do~danego wierzchołka musi być równa sumie wartości przepływu na~krawędziach
wychodzących z~tego wierzchołka.

\begin{defn}
\textbf{Wartością przepływu} $f$ w~sieci $S$ nazywamy liczbę
$$ W(f) = \sum_{u: \; su \in E} f(su) - \sum_{u: \; us \in E} f(us) $$
\end{defn}

\noindent
Powyższe definicje wystarczą, aby~zdefiniować problem maksymalnego przepływu.

\begin{defn}[Problem maksymalnego przepływu]
Dana jest sieć $S = (G,c,s,t)$. Szukamy przepływu $f$ o~maksymalnej wartości
$W(f)$, zwanego również \textbf{przepływem maksymalnym}.
\end{defn}

\noindent
Zagadnienie znajdowania maksymalnego przepływu jest rozwiązywalne przez
wiele zachłannych algorytmów opartych na metodzie Forda-Fulkersona,
polegającej na~znajdowaniu ścieżek w~tzw. sieci rezydualnej. Szczegółowy
opis jednego z~takich algorytmów znajduje się w następnej sekcji.

\section{Algorytm}
\label{sec:algorithm}

\section{Dowód poprawności}
\label{sec:correctnessproof}

W~tej sekcji wykażemy związek między postawionym problemem a~zagadnieniem
wyznaczania przepływu maksymalnego oraz równoważność rozwiązań obu zadań.\\

\noindent
Na początek zdefiniujmy w~sposób formalny pojęcia użyte w~oryginalnym zadaniu.
Załóżmy, że dane są następujące zbiory:

\begin{itemize}
	\item \textbf{zbiór ekspertów}, oznaczony $E$,
	\item \textbf{zbiór umiejętności}, oznaczony $S$,
	\item \textbf{zbiór projektów}, oznaczony $P$.
\end{itemize}

\begin{defn}
\textbf{Funkcją umiejętności} nazywamy funkcję
$$ \ability : E \times S \to \{ 0,1 \} $$
gdzie dla eksperta $e \in E$ oraz umiejętności $s \in S$ zachodzi
$\ability(e, s) = 1$ wtedy i~tylko wtedy, gdy ekspert $e$ posiada umiejętność
$s$, zaś 0 w przeciwnym przypadku.
\end{defn}

\begin{defn}
\textbf{Zapotrzebowaniem projektu} nazywamy funkcję
$$ \need : P \times S \to \mathbb{N} $$
gdzie dla projektu $p \in P$ i~umiejętności $s \in S$ zachodzi $\need(p, s) = n$
wtedy i~tylko wtedy, gdy w~projekcie $p$ liczba potrzebnych ekspertów
w~dziedzinie umiejętności $s$ wynosi $n$.
\end{defn}

\begin{defn}
\textbf{Przyporządkowaniem eksperta} nazywamy funkcję
$$ \assign : E \to P \times S $$
gdzie dla projektu $p \in P$ i~umiejętności $s \in S$ zachodzi
$\assign(e) = (p,s)$ wtedy i~tylko wtedy, gdy:
\begin{itemize}
	\item ekspert $e$ posiada umiejętność $s$ (tj. $\ability(e,s) = 1$),
	\item ekspert $e$ został przyporządkowany do pracy w~projekcie $p$
	w~dziedzinie umiejętności $s$.
\end{itemize}
Ponadto, dla każdego projektu $p$ i umiejętności $s$ musi zachodzić
$$ \left| \assign^{-1}[p,s] \right| \leq \need(p,s) $$
gdzie $\assign^{-1}[p,s]$ jest przeciwobrazem funkcji $\assign$ dla argumentów
$p,s$.
\end{defn}

\begin{defn}
\textbf{Liczbą braków w projekcie $p$} dla danego przyporządkowania $\assign$
nazywamy liczbę
$$ \missing(p,\assign) = \sum_{s \in S} \left( \need(p,s) -
\left| \assign^{-1}[p,s] \right| \right) $$
\end{defn}

\begin{defn}
\textbf{Całkowitą liczbą braków} dla danego przyporządkowania $\assign$ nazywamy
liczbę
$$ M(\assign) = \sum_{p \in P} \missing(p, \assign) $$
\end{defn}

\noindent
Widoczne jest, że $M$ jest parametrem minimalizowanym w~postawionym problemie,
zależnym od~końcowego przyporządkowania.\\

\noindent
Na~podstawie powyższych definicji skonstruujemy teraz sieć, której użyjemy
do~wyznaczenia rozwiązań problemu.

\begin{defn}
\textbf{Siecią przydziałów} nazwiemy sieć $N = \left(G,c,s_G,t_G\right)$, gdzie:
\begin{itemize}
	\item $G = \left(V_G,E_G\right)$ jest grafem skierowanym takim, że:
	\begin{itemize}
		\item $V_G = E \cup S \cup P \cup \left\lbrace s_G,t_G \right\rbrace$,
		\item $E_G = \left\lbrace (e,s) : \ability(e,s) = 1, e \in E, s \in S
		\right\rbrace \cup \left\lbrace (s,p) : \need(s,p) > 0, s \in S, p \in P
		\right\rbrace$, tj. krawędziami połączeni są eksperci z~ich opanowanymi
		umiejętnościami, oraz projekty z~potrzebnymi do~ich realizacji
		umiejętnościami.
	\end{itemize}
	\item $c : E_G \to \mathbb{N}$ jest funkcją pojemności zdefiniowaną
	następująco:
	\begin{itemize}
		\item jeżeli $e = s_G u, u \in E$, to $c(e) = 1$,
		\item jeżeli $e = es, e \in E, s \in S$, to $c(e) = \ability(e,s) = 1$,
		\item jeżeli $e = sp, s \in S, p \in P$, to $c(e) = \need(p,s)$,
		\item jeżeli $e = pt_G, p \in P$, to
		$$c(e) = \sum_{sp \in E_G} c(up)$$
		(tj. pojemność tej krawędzi jest równa sumie pojemności krawędzi
		wchodzących do wierzchołka $p$).
	\end{itemize}
	\item $s_G,t_G$ są wyróżnionymi wierzchołkami z $V_G$ --- kolejno źródłem
	i~ujściem.
\end{itemize}
\end{defn}

\begin{lem}
Dla dowolnego przyporządkowania $\assign$ mamy
$$ M(\assign) \geq \left( \sum_{p \in P} \sum_{s \in S} \need(p,s) \right)
- |E| $$
\end{lem}

\begin{proof}
Mamy
\begin{align*}
M(\assign) &= \sum_{p \in P} \missing(p, \assign) = \\
&= \sum_{p \in P} \left( \sum_{s \in S} \need(p,s) -
\left| \assign^{-1}[p,s] \right| \right) = \\
&= \left( \sum_{p \in P} \sum_{s \in S} \need(p,s) \right) -
\left(\sum_{p \in P} \sum_{s \in S} \left| \assign^{-1}[p,s] \right|\right) = \\
&\geq \left( \sum_{p \in P} \sum_{s \in S} \need(p,s) \right) - |E|
\end{align*}
na mocy rozłączności przeciwobrazów funkcji $\assign$.
\end{proof}

\begin{thm}
Przepływ maksymalny w~sieci przydziałów wyznacza przyporządkowanie o~minimalnej
możliwej wartości parametru $M$. 
\end{thm}

\begin{proof}
Aby dowieść to twierdzenie, wykażemy kolejno, że:
\begin{enumerate}
	\item Dowolny przepływ w~sieci przydziałów wyznacza poprawne (niekoniecznie
	maksymalne) przyporządkowanie ekspertów do~zadań.
	
	Niech dany będzie pewien przepływ $f$ w~sieci przydziałów $N$.
	Przyporządkowanie ekspertów do projektów wyznaczamy w~następujący sposób:
	\begin{enumerate}
		\item Pewnego eksperta $e \in E$ przypisujemy do~umiejętności $s \in S$,
		jeżeli $f(e,s) = 1$.
		%\item Pewnemu projektowi $p \in P$ w~dziedzinie umiejętności $s \in S$
		%przypisujemy dokładnie $f(s,p)$ ekspertów.
		\item Niech dana będzie pewna umiejętność $s \in S$. Oznaczmy zbiór
		ekspertów przypisanych do~tej umiejętności w punkcie (a) jako $E_s$.

		Zbiór $E_s$ dzielimy na~rozłączne podzbiory $E_{s,p}$ takie, że
		$|E_{s,p}| = f(s,p)$.
		\item Dla każdego z uzyskanych podzbiorów $E_{s,p}$, gdzie
		$s \in S, p \in P$ określamy wartość przyporządkowania $\assign_f$ jako
		$$ (\forall e \in E_{s,p}) \assign_f(e) = (s,p) $$

	\end{enumerate}
	Zauważmy następujące fakty:
	\begin{itemize}
		\item Rozważmy wierzchołek $e \in E$ odpowiadający pewnemu ekspertowi.
		
		Z~definicji zbioru krawędzi sieci i~funkcji przepustowości,
		do~wierzchołka tego wchodzi dokładnie jedna krawędź o~pojemności 1,
		a~wychodzi z~niego co~najwyżej $|S|$ krawędzi o~pojemności 1.

		Stąd w~przepływie $f$ tylko jedna z~krawędzi wychodzących może mieć
		przepływ 1, a~więc każdy ekspert może być przyporządkowany
		do~co~najwyżej jednej umiejętności.
		\item Rozważmy dowolny wierzchołek $s \in S$ odpowiadający pewnej
		umiejętności.

		Z~własności przepływu mamy
		$$ \sum_{us \in E_G} f(us) = \sum_{sv \in E_G} f(sv) $$
		Wiedząc, że~wszystkie krawędzie wchodzące do~$s$ wychodzą ze~zbioru $E$,
		oraz że~wszystkie krawędzie wychodzące z~$s$ wchodzą do~zbioru $P$, mamy
		$$ \sum_{e \in E} f(es) = \sum_{p \in P} f(sp) $$
		Krawędzie o~niezerowym przepływie wchodzące do~$e$ reprezentują
		ekspertów przydzielonych do~danej umiejętności, zaś krawędzie
		o~niezerowym przepływie wychodzące z~$e$ reprezentują zapotrzebowanie
		projektów na~ekspertów z~umiejętnością $s$.

		Ponieważ suma przepływów krawędzi wchodzących i~wychodzących jest
		taka sama, każdego eksperta w~dziedzinie umiejętności $s$
		można przypisać do~dokładnie jednego projektu, a~więc można wykonać
		punkt (b) konstrukcji.
		\item Rozważmy dowolne dwa wierzchołki $s \in S, p \in P$ takie, że
		$sp \in E_G$. Z~definicji sieci mamy $c(s,p) = \need(s,p)$,
		a~z~konstrukcji rozwiązania wynika, że
		$f(s,p) = \left|\assign_f^{-1}[s,p]\right|$. Stąd na~mocy definicji
		przepływu mamy
		$$ \left|\assign_f^{-1}[s,p]\right| = f(s,p) \leq c(s,p) = \need(s,p)$$
		\item Rozważmy dowolny wierzchołek $p \in P$. Z~definicji funkcji
		pojemności, jeśli wszystkie krawędzie wchodzące do~$p$ będą wysycone
		przepływem (tj. $f(e) = c(e)$), to~przepływ ten można przekazać
		w~całości do ujścia krawędzią $pt_G$, bo
		$$ c(pt_G) = \sum_{up \in E_G} c(up) $$
		Stąd pojemność krawędzi $pt_G$ nie ogranicza wartości maksymalnego
		przepływu.
	\end{itemize}
	Wyznaczone przyporządkowanie $\assign_f$ spełnia więc wszystkie warunki
	prawidłowego przyporządkowania ekspertów do projektów.

	\item Skonstruowane z~przepływu maksymalnego przyporządkowanie wyznacza
	minimalną wartość parametru $M$.

	Dowód nie~wprost: Załóżmy, że~przyporządkowanie $\assign_{\max}$
	skonstruowane na~podstawie przepływu maksymalnego $f_{\max}$ w~sposób
	opisany w punkcie 1. można rozszerzyć tak, aby zmniejszyć wartość parametru
	$M$.

	Na mocy lematu mamy
	\begin{align*}
	\left( \sum_{p \in P} \sum_{s \in S} \need(p,s) \right) -
	\left(\sum_{p \in P} \sum_{s \in S} \left| \assign_{\max}^{-1} [p,s]\right|\right) 
	&= M(\assign_{\max}) > \\
	&> \left( \sum_{p \in P} \sum_{s \in S} \need(p,s) \right) - |E|
	\end{align*}
	a~więc istnieje co~najmniej jeden ekspert $e \in E$, który nie~jest
	przyporządkowany do~jednego z~projektów.
	\begin{enumerate}
		\item Jeżeli $(\forall s \in S) \ability(e,s) = 0$, to eksperta
		nie~można przyporządkować do~żadnego projektu.
		\item Załóżmy, że istnieje umiejętność $s \in S$ taka, że
		$\ability(e,s) = 1$.
		\begin{enumerate}
			\item Jeżeli $(\forall p \in P) \need(p,s) =
			\left|\assign_{\max}^{-1} [p,s]\right|$, to eksperta nie~można
			przyporządkować do~żadnego projektu.
			\item Załóżmy, że~istnieje projekt $p$ taki, że
			$\need(p,s) > \left|\assign_{\max}^{-1} [p,s]\right|$.
			Wówczas w~sieci rezydualnej dla przepływu $f_{\max}$ istnieje
			ścieżka rozszerzająca postaci $s_G e s p t_G$, wzdłuż której można
			zwiększyć przepływ o~1. Sprzeczność z~maksymalnością przepływu
			$f_{\max}$.
		\end{enumerate}
	\end{enumerate}
\end{enumerate}
\end{proof}

\end{document}
